\section{Naturalité conjointe des publications de l'état généalogique}
\label{sec:cpe-naturalidad-tres-publicaciones}

Soit \((X_m,p_{m+1,m})_{m\ge0}\) le système projectif d'états
APP--TRIT--TPK étalonnés. À chaque niveau sont conservés la fibre, la relation,
le nombre, l'orientation, la retenue, la mémoire, la frontière, la route, la multisection,
la survie, le terminal et le lecteur. Soient
\[
 \mathfrak R_m:X_m\to A_m,
 \qquad
 \mathfrak E_{m,\chi}:X_m\to E_m,
 \qquad
 \mathfrak D_m:X_m\to D_m
\]
les publications archimédienne, exceptionnelle étalonnée et dimensionnelle. Leurs
troncatures sont notées \(a_{m+1,m}\), \(e_{m+1,m}\) et
\(d_{m+1,m}\).

\begin{theorem}[Cône naturel des \(3\) publications]
\label{thm:cpe-cono-tres-publicaciones}
Si les lecteurs de niveau fini satisfont
\begin{equation}
\begin{aligned}
 a_{m+1,m}\mathfrak R_{m+1}&=\mathfrak R_mp_{m+1,m},\\
 e_{m+1,m}\mathfrak E_{m+1,\chi}&=\mathfrak E_{m,\chi}p_{m+1,m},\\
 d_{m+1,m}\mathfrak D_{m+1}&=\mathfrak D_mp_{m+1,m},
\end{aligned}
\label{eq:cpe-naturalidad-tres-lectores}
\end{equation}
alors les lecteurs de limite existent et sont uniques
\[
 \mathfrak R_\infty:X_\infty\to A_\infty,
 \quad
 \mathfrak E_{\infty,\chi}:X_\infty\to E_\infty,
 \quad
 \mathfrak D_\infty:X_\infty\to D_\infty
\]
qui reproduisent tous les niveaux. L'application conjointe
\begin{equation}
 \mathfrak J_\chi
 :=(\mathfrak R_\infty,\mathfrak E_{\infty,\chi},\mathfrak D_\infty)
 :X_\infty\longrightarrow A_\infty\times E_\infty\times D_\infty
 \label{eq:cpe-mapa-conjunto-tres-publicaciones}
\end{equation}
est une publication du même état et non une concaténation de trois
générateurs.
\end{theorem}

\begin{proof}
Pour \(x=(x_m)_m\in X_\infty\), les trois suites
\((\mathfrak R_mx_m)_m\), \((\mathfrak E_{m,\chi}x_m)_m\) et
\((\mathfrak D_mx_m)_m\) sont compatibles par
\eqref{eq:cpe-naturalidad-tres-lectores} ; elles définissent, par la propriété universelle
de la limite projective, les trois lecteurs. La même propriété prouve leur unicité.
La formule \eqref{eq:cpe-mapa-conjunto-tres-publicaciones} évalue une seule
histoire \(x\) et conserve donc l'antécédent causal commun.
\end{proof}

L'égalité légitime n'affirme pas qu'une sortie avec perte reconstruit à elle
seule toutes les autres. Elle affirme que les trois carrés de troncature commutent et
que leurs sorties sont des coordonnées naturelles d'un même cône. En particulier,
la réalisation archimédienne peut oublier la généalogie sans que cet oubli soit
attribué à la publication exceptionnelle ou à la publication dimensionnelle.

\section{Séparation conjointe de la valeur et de l'incidence exceptionnelle}
\label{sec:cpe-separacion-conjunta}

La jauge \(\chi_{\rm exc}\) étant fixée, le lecteur exceptionnel agit sur chaque
générateur par
\begin{equation}
 (w_j,f_j)\longmapsto
 \bigl(f_j,\Gamma_{f_jA_Wf_j^{-1}}(w_j)\bigr).
 \label{eq:cpe-lector-excepcional-generador}
\end{equation}
Le cadre \(f_j\) fait partie de la sortie et
\(\Gamma_{f_jA_Wf_j^{-1}}\) est un automorphisme de la fibre étalonnée.

\begin{theorem}[Séparation étalonnée des histoires]
\label{thm:cpe-separacion-calibrada}
Le lecteur exceptionnel
\(\mathfrak E_{\infty,\chi_{\rm exc}}\) est injectif sur le quotient par
les changements de coordonnées inclus dans \(\chi_{\rm exc}\). Par conséquent,
\[
 (\mathfrak R_\infty,\mathfrak E_{\infty,\chi_{\rm exc}})
 :X_\infty^{\rm cal}\longrightarrow
 A_\infty\times E_\infty
\]
est un plongement topologique lorsque \(X_\infty^{\rm cal}\) est compact et que le
codomaine est séparé.
\end{theorem}

\begin{proof}
À partir de la sortie \((f_j,y_j)\), où
\(y_j=\Gamma_{f_jA_Wf_j^{-1}}(w_j)\), on retrouve
\[
 w_j=\Gamma_{f_jA_Wf_j^{-1}}^{-1}(y_j).
\]
Cela reconstruit chaque générateur et, par compatibilité de troncature, toute l'
histoire. Deux représentants qui ne diffèrent que par un changement inclus dans la
jauge définissent la même classe ; en dehors de cette équivalence, la reconstruction
est unique. L'injectivité conjointe est immédiate. Une application continue et
injective d'un compact vers un espace séparé est un homéomorphisme sur son image.
\end{proof}

\section{Antécédent non commutatif commun de la double projection}
\label{sec:cpe-antecedente-ucp}

Soit
\(\mathcal A_{\rm TPK}=C(X_F^{\leftrightarrow})\rtimes_\sigma\mathbb Z\)
et soit \(u\) l'unitaire qui implémente \(\sigma\). Les lecteurs continus
\(\mathfrak R\) et \(\mathfrak E_\chi\) induisent les morphismes unitaux
\begin{equation}
 \Phi_{\rm arq}(f)=f\circ\mathfrak R,
 \qquad
 \Phi_{\rm exc}(g)=g\circ\mathfrak E_\chi,
 \label{eq:cpe-dos-mapas-ucp}
\end{equation}
des algèbres commutatives de leurs codomaines vers
\(C(X_F^{\leftrightarrow})\subset\mathcal A_{\rm TPK}\).

\begin{theorem}[Système d'opérateurs commun et équivariance TPK]
\label{thm:cpe-sistema-operadores-comun}
Le sous-espace autoadjoint unital
\[
 \mathcal S_\infty
 =\operatorname{span}\bigl\{
 I,u,u^*,\Phi_{\rm arq}(C(A_\infty)),
 \Phi_{\rm exc}(C(E_\infty))\bigr\}
 \subset\mathcal A_{\rm TPK}
\]
est un système d'opérateurs non commutatif. Les deux applications de
\eqref{eq:cpe-dos-mapas-ucp} sont UCP\@. Si les lecteurs entrelacent les
dynamiques de sortie \(\alpha_{\rm arq}\) et \(\alpha_{\rm exc}\), alors
\begin{equation}
 \operatorname{Ad}_u\Phi_{\rm arq}
 =\Phi_{\rm arq}\alpha_{\rm arq},
 \qquad
 \operatorname{Ad}_u\Phi_{\rm exc}
 =\Phi_{\rm exc}\alpha_{\rm exc}.
 \label{eq:cpe-equivarianza-doble-ucp}
\end{equation}
\end{theorem}

\begin{proof}
La composition avec une application continue est un homomorphisme unital de
\(C^*\)-algèbres et est donc complètement positive. L'unitaire
\(u\) rend l'enveloppe non commutative lorsque la dynamique est non triviale.
L'identité \eqref{eq:cpe-equivarianza-doble-ucp} se vérifie en évaluant
les deux membres en une histoire et en utilisant l'équivariance de chaque lecteur.
\end{proof}

Cette construction ne force pas deux mesures qutrit mutuellement non biaisées à
posséder une PVM conjointe. L'antécédent commun est l'algèbre croisée des
histoires ; les publications sont deux canaux UCP distincts depuis leurs
codomaines.

